%% ============================================================
%% 分片 35：原文第 12 节「Branes From Gauge Theory」（原文 glshort.tex 起于第 11186 行）
%% 由 tools/harvmac2tex.py 从 arXiv 源码（hep-th/0604151v3）自动生成。
%% 只译英文散文；公式、编号、\label、\cite 一律不改。
%% ============================================================

\section{规范场论中的膜}
\label{sec:zapendix}

\begin{keytakeaways}
\begin{itemize}[leftmargin=1.4em,itemsep=2pt,topsep=2pt]
\item \textbf{物理图像}：本文关心的几类膜——零膜（$(B,B,B)$）、Hitchin 纤维（$(B,A,A)$）、典范余各向同性膜（$(A,B,A)$）——都可以提升为四维 ${\EUN}=4$ 超 Yang-Mills 理论中的边界条件。边界条件保持超对称的判据是超流的法向分量在边界上为零；由于约化到二维并扭化会再砍掉一半超荷，规范场论中保持四分之一超对称的膜，在 sigma 模型语言里称作 $1/2$ BPS 膜。
\item \textbf{数学工具}：用十个 $\Gamma_I$ 与 $\Gamma_{2367}$ 把超荷投影到扭化后保留下来的子空间，边界条件写成 $T\epsilon=\epsilon$、$\Gamma_{2367}\epsilon=\epsilon$；由费米子的边条件反推玻色子的边条件，再验证 $\bar\epsilon J_1=0$ 与 $\delta\lambda$、$\delta A_I$ 的变分相容；边界值上的 Hitchin 方程由 $F_{IJ}$ 在特定子空间内自对偶而自然出现。
\item \textbf{关键结论}：取 $T=\Gamma_{04}$ 并要求 $T\lambda\vert=\lambda\vert$，会推出 $F_{0I}+F_{4I}=0$（$I\neq0,4$）与 Hitchin 方程 $F_{23}=F_{67}$ 等，于是规范场在边界上的取值在 $\MH$ 中挑出一点 $p$，得到支撑在 $p$ 的零膜。改用 $\tilde T=T\Gamma_{2367}$ 可得另一构造：它在平坦时空中只保四个超荷，但在 $\MH$ 的光滑点处与旧构造在低能理论中等价。
\item \textbf{边界与展望}：当 $p$ 对应可约 Higgs 丛时，两种零膜构造未必等价（$A_4,A_5,A_8,A_9$ 的零模给出 Higgs 丛的对称性）；另一半超对称边界条件只在远离奇点处才与二维图像相符。事实上提升到规范场论并不唯一，本文只给出足以展示关键膜的少数简单构造。
\end{itemize}
\end{keytakeaways}

至此我们一直在靶为 $\MH$ 的有效二维 sigma 模型中考虑膜。本结论节的目标是
说明：重要的膜其实可以在四维规范场论中定义。

任何有价值的膜总是至少保持一个超对称。例如，它们是复结构 $J$ 中的
$B$-膜或复结构 $K$ 中的 $A$-膜。但对几何 Langlands 对偶重要的膜，关于
三个复结构都有特殊的性质。例如，$\MH(\LG,C)$ 上的零膜是一个
$(B,B,B)$-膜：它关于复结构 $I,J,K$ 都是 $B$-膜。它的镜像对象是
Hitchin 纤维化的一条纤维，该纤维关于 $I$ 是全纯的、关于 $J$ 与 $K$ 是
Lagrangian 的。因此它是 $(B,A,A)$-膜。典范余各向同性膜是
$(A,B,A)$-膜，它的镜像对象也是。这些类型中任何一种膜都保持两个拓扑
超荷（它们的线性组合例如对 $(B,B,B)$ 型膜给出复结构 $I$、$J$、$K$ 中
所需的 $B$-型超对称）。

我们将集中讨论这四类膜，并描述规范场论中相应的边界条件。提升到规范
场论并不一定唯一；事实上，在某些情形我们会描述规范场论中若干不同的
边界条件，它们在约化到二维后显然变得等价——至少远离奇点时如此。

为把扭化后的理论讨论为拓扑场论，人们把四维 Minkowski 空间换成一般的
四维流形 $M$。在拓扑场论中，最方便的是取 $M$ 具有欧氏号差，理由如下。
走向一般四维流形时，为保持一些超对称，人们通过把前四个标量场
$\phi_0,\dots,\phi_3$ 解释为 $M$ 的切丛的一个截面来做「扭化」。由于标量
场天然具有正号差，若 $M$ 取欧氏号差，这一扭化就自然得多。（另一种做法
也许是在 $\phi$ 空间做 Wick 转动，但这似乎并不妥当。）此外，任何紧四维
流形都容许正号差的度量，而容许 Lorentz 号差度量则是严重的拓扑限制。

就目前的讨论而言，我们将在更受限的情形考虑膜。照例，我们在黎曼面 $C$
上紧化到二维。这给出一个有效二维理论，而我们要理解这个理论中的膜如何
能在底层的四维规范场论中得到解释。在边界附近存在一个自然的「时间」
方向（垂直于 $C$、切于边界），做 Wick 转动并无太大好处。相反，我们愿意
避免 Wick 转动，以便清楚地表明我们所考虑的一切膜都是物理上合理且酉的。

既然如此，只要拓扑超荷 $Q$ 被我们的边界条件保持，其伴随 $Q^\dagger$ 也
被保持。它们服从物理的超对称代数 $\{Q,Q^\dagger\}=H$，其中 $H$ 是
哈密顿量。由于我们的边界条件保持两个拓扑超荷，它们共保持四个超对称。
这是 ${\EUN}=4$ 超 Yang-Mills 理论 16 个整体超对称的四分之一，因此这类
膜称为 $1/4$ BPS 膜。从 sigma 模型（其超对称只有一半）的角度看，它们
称为 $1/2$ BPS 膜。然而回到规范场论，在某些情形我们在规范场论中的边界
条件沿 $C$ 局部保持八个超对称，而约化为四个超对称来自 $C$ 的扭化与曲率。

为描述二维有效理论中的膜，我们在（比如）$\Bbb{R}^2_+$ 上表述该理论，
坐标为 $x^0$、$x^1$，其中 $-\infty< x^0< \infty$，$0\leq x^1 <\infty$。
超对称膜通过在 $x^1=0$ 上指定一个超对称边界条件而得到。要从底层四维
规范场论得到这样的构造，我们只需在 $M=\Bbb{R}^2_+\times C$ 上考虑该
理论，并同样在 $x^1=0$ 处描述超对称边界条件。

我们用 $x^2,x^3$ 表示 $C$ 上的局部坐标。由于取 $\Bbb{R}^2_+$ 为平坦的，
扭化只影响两个标量场 $\phi_2,\phi_3$（或 $A_6$, $A_7$；记号已在
第~\ref{sec:reviewsym}~节解释过）。紧化与扭化未破缺的超对称由那些在
适当意义下于 $C$ 上协变常数的旋量 $\epsilon$ 生成。在扭化后的理论中，
条件简单地就是 $\epsilon$ 应满足 % 式 (12.1)
\setcounter{equation}{0}
\begin{equation}\label{eq:ucuv}
\Gamma_{2367}\epsilon=\epsilon,
\end{equation}
 正如式 \eqref{eq:ytytag} 那样。

\subsection{边界条件的一般性质}

首先让我们回忆关于边界条件应当是什么的若干一般性事实。

边界条件以这样一种方式约束场的值（对服从二阶运动方程的玻色子，还约束
其法向导数），使得运动方程中的边界项为零。例如，对具有号差
$-++\dots+$ 的度量 $g$ 的 $n$ 维流形 $M$ 上的规范场，最小作用量为 % 式 (12.2)
\setcounter{equation}{1}
\begin{equation}\label{eq:upco}
{{\cmmib I}}_B={1\over 2e^2}\int_M
d^nx \sqrt g\,\Tr \,F_{IJ}F^{IJ}.
\end{equation}
 $M$ 是 $n$ 维流形，以 $x^0$ 为「时间」方向，边界由 $x^1=0$ 定义。
作用量 $\delta {{\cmmib I}}_B$ 的变分中的边界项是 % 式 (12.3)
\setcounter{equation}{2}
\begin{equation}\label{eq:ofof}
(\delta {\cmmib
I}_B)_{bdry}={2\over e^2}\int_{\partial M}d^{n-1}x \sqrt
g\,\Tr\left(g^{IJ}\sum_{I,J\not=1}\delta A_I F_{1J}\right).
\end{equation}
 边界条件必须把 $F_{I1}$ 与 $\delta A_I$（$I\not= 1$）的边界值的某个
线性组合置零，使得 $(\delta {{\cmmib I}}_B)_{bdry}=0$。对某个 $I$ 置
$\delta A_I=0$ 的含义是：指定 $A_I$ 的边界值，换言之，在 $A_I$ 上施加
Dirichlet 边界条件（边界值取定）。反过来，在 $M$ 边界附近取 $A_1=0$ 的
规范时，条件 $F_{I1}=0$ 意味着 $A_I$ 的法向导数为零。所以 $F_{I1}=0$
相当于 $A_I$ 上 Neumann 边界条件的规范不变版本。为便于量子化，我们
不想过度约束边界值，因此把 $\delta A_I$ 与 $F_{I1}$ 的一个最小线性
组合置零，使得 $(\delta {{\cmmib I}}_B)_{bdry}=0$。（因此边界条件在
适当意义上定义了边界数据的一个 Lagrangian 子流形。）也还可以给
${{\cmmib I}}_B$ 添加边界项，从而修改 $(\delta {{\cmmib I}}_B)_{bdry}$；
边界条件总是定义允许的边界值（$A_I$ 与 $F_{I1}$，$I\not= 1$）的最大
子空间，使得 $(\delta {{\cmmib I}}_B)_{bdry}=0$。我们把这样的边界条件
称为双曲边界条件。

费米子当然也有类似的故事。最小作用量为 % 式 (12.4)
\setcounter{equation}{3}
\begin{equation}\label{eq:olpo}
{{\cmmib I}}_F=-{i\over e^2}\int_M d^nx \sqrt
g\,\Tr\,{\bar\lambda} \Gamma^I D_I\lambda,
\end{equation}
 人们可以给它添加额外的边界项。用最小作用量时，$\delta {{\cmmib
I}}_F$ 中的边界项是 % 式 (12.5)
\setcounter{equation}{4}
\begin{equation}\label{eq:rolpo}
\delta
{{\cmmib I}}_F=-{i\over e^2}\int_{\partial M}d^{n-1}x\sqrt
g\,\Tr\,{\bar \lambda} \Gamma_1\delta\lambda.
\end{equation}
 对费米子，合适的边界条件把 $\lambda$ 的一半分量的边界值（从而
$\delta\lambda$ 的一半分量）置零，使得 ${\bar \lambda}
\Gamma_1\delta\lambda=0$。

在体理论中，超对称意味着（当 $M$ 平坦时，或当 $M$ 弯曲而做了适当扭化
后）存在守恒的超流 $J_I$。对超对称 Yang-Mills 理论，
% 式 (12.6)
\setcounter{equation}{5}
\begin{equation}\label{eq:noggin}
J_I={1\over 2}\G^{KL}F_{KL}\G_I\lambda.
\end{equation}
 在存在边界时，边界条件与超对称相容的条件是 $J$ 的法向分量在边界上
为零，从而超流的通量不会在边界处消失。更精确地说，条件是 % 式 (12.7)
\setcounter{equation}{6}
\begin{equation}\label{eq:hobbo}
\bar \epsilon J_1=0
\end{equation}
 对适当的超对称生成元 $\epsilon$。

这自动保证了边界条件与玻色子、费米子的超对称变分都相容。因此，如果
已经知道要在玻色子和费米子上取什么边界条件，$J_1=0$ 就是基本的条件。
实际上，由于一开始并不知道想要什么边界条件，我们将先假定费米子上有
一个简单的边界条件（选得使 ${\bar \lambda}\Gamma_1\delta\lambda=0$），
再推出玻色子上相应的边界条件以保证费米子的边界条件在超对称下保持，
然后验证 $J_1=0$。

若 $J_1=0$，我们就找到了能够定义守恒超荷的边界条件（当 $M$ 平坦时；
当 $M$ 弯曲时，则在适当扭化的理论中）。然而我们还想要更多。如果目标是
研究超对称态或拓扑场论（其中只关心超对称态），那么迫使超对称自发破缺
的边界条件就不太有意义。玻色子的 Dirichlet 边界条件一般与未破缺的
超对称不相容。如我们所讨论的，Dirichlet 边界条件意味着规范场的某些
分量 $A_I$ 的边界值、从而（可能）场强的某些分量 $F_{IJ}$ 的边界值被
指定。我们主要关心这样的边界条件：$F_{IJ}$ 的边界值以一种与未破缺
超对称相容的方式被固定，即要求超对称变分
% 式 (12.8)
\setcounter{equation}{7}
\begin{equation}\label{eq:undop}
\delta\lambda =\half\G^{IJ}F_{IJ}\epsilon
\end{equation}
 对适当的超对称生成元 $\epsilon$ 为零。

\subsection{$(B,B,B)$ 型膜}

我们要构造的膜总是保持那些生成元满足适当条件 $T\epsilon=\epsilon$
以及\footnote{把条件 $T\epsilon=\epsilon$ 换成
$T\epsilon=-\epsilon$ 不会有本质上的新东西，因为理论中总存在使 $T$
变号的对称性。例如在我们的第一个例子里 $T=\Gamma_{04}$，这样的对称性
就是 $49$ 平面内的反射。} $\Gamma_{2367}\epsilon=\epsilon$ 的超对称。
两个条件合起来将允许四个超对称。$T$ 将与 $\Gamma_{2367}$ 交换。在某些
例子中，$T$ 的选取使 $T^2=1$。此时 $T\epsilon=\epsilon$ 这一条件本身
允许八个超对称，而边界条件将保持这八个超对称。即便如此，紧化到二维
以及「扭化」会把第二个条件 $\Gamma_{2367}\epsilon=\epsilon$ 加在未破缺
的超对称上，在低能理论中只留下四个超对称。

我们将说明一些能给出各类膜的 $T$ 的示例性选择。我们当然不声称描述了
四维规范场论中所有有意思的膜的构造，只描述若干简单的构造，它们足以
在规范场论中展示本文最重要的那些膜。

要得到 $(B,B,B)$ 型膜，我们先取 $T=\G_{04}$ 并施加边界条件 % 式 (12.9)
\setcounter{equation}{8}
\begin{equation}\label{eq:unvu}
T\lambda\vert
=\lambda\vert.
\end{equation}
 对任何场 $\Phi$，我们用 $\Phi|$ 表示 $\Phi$ 在 $x^1=0$ 上的限制。
（不过，当不太可能引起混淆时，我们省略这个符号。）$\lambda$ 上的这个
边界条件只与这样的超对称相容：超对称变换
$\lambda\to \lambda+\delta\lambda$ 保持条件
$T\lambda\vert=\lambda\vert$。要求就是 $(T-1)\delta\lambda=0$。由于
$\delta\lambda=\half\G^{IJ}F_{IJ}\epsilon$ 且
$T\epsilon=\epsilon$，我们有 % 式 (12.10)
\setcounter{equation}{9}
\begin{equation}\label{eq:bogus}
(T-1)\delta\lambda =\half [T,\G^{IJ}F_{IJ}]\epsilon.
\end{equation}

计算对易子后，我们发现需要
% 式 (12.11)
\setcounter{equation}{10}
\begin{equation}\label{eq:nogus}
\sum_{I\not=0,4}(\Gamma_{0I}F_{4I}-\Gamma_{4I}F_{0I})\epsilon=0.
\end{equation}

由于 $\epsilon$ 恰好受 $\Gamma_0\Gamma_4\epsilon=\epsilon$ 约束，这一
条件等价于 % 式 (12.12)
\setcounter{equation}{11}
\begin{equation}\label{eq:hogus}
F_{0I}+F_{4I}=0~~{\rm for}~I\not=0,4.
\end{equation}

为选取与上述描述相容的双曲边界条件，我们（在适当的规范下）把
$A_0+A_4$ 的边界值置零，而对 $A_I$（$I\not= 0,1,4$）施加 Dirichlet
边界条件，其边界值取定且不依赖时间。最后，\eqref{eq:hogus}\ 告诉我们
$F_{01}+F_{41}=0$，这对应于 $A_0+A_4$ 上规范不变的 Neumann 边界条件。
（由于 $A_0+A_4$ 对应一个「零」方向，在双曲边界条件中可以同时把
$A_0+A_4$ 及其法向导数置零。这依赖于号差中的负号，即
\eqref{eq:ofof} 中 $g^{IJ}={\rm diag}(-1,1,1,\dots,1)$ 这一事实。）

到目前为止，我们从费米子上的边界条件出发，猜测了玻色子上应有的边界
条件。在这个阶段，我们可以验证基本条件：超流的法向分量在边界上为零：
% 式 (12.13)
\setcounter{equation}{12}
\begin{equation}\label{eq:jolygo}
{\bar\epsilon}\G^{MN}F_{MN}\G_1\lambda=0.
\end{equation}
 由于 $\G_{04}\lambda=\lambda$ 且 $\G_{04}\epsilon=\eps$，只有那些
$F_{MN}$ 中恰有一个指标等于 $0$ 或 $4$ 的分量才有贡献；利用这一点可以
验证：若 $F_{0I}+F_{4I}=0$（$I\not= 0,4$），则 \eqref{eq:jolygo}\ 成立。

这意味着费米子与玻色子上的边界条件都与超对称相容。对费米子我们已
看到这一点，对玻色子也可直接验证。玻色子的超对称变分
% 式 (12.14)
\setcounter{equation}{13}
\begin{equation}\label{eq:guffo}
\delta A_I=i\bar\epsilon\Gamma_I\lambda
\end{equation}
 连同 $(\Gamma_{04}-1)\lambda=(\G_{04}-1)\epsilon=0$ 蕴含：对
$I\not= 0,4$ 有 $\delta A_I=0$，且 $\delta A_0+\delta A_4=0$，从而
$\delta(F_{0I}+F_{4I})=0$（$I\not=0,4$）。因此我们的条件
$F_{0I}+F_{4I}=0$（$I\not= 0,4$）与超对称相容。

到此为止，我们得到了一个保证守恒超荷存在的边界条件。我们还希望边界
条件容许可存在超对称态。为此，$A_I$（$I\not=0,1,4$）上的 Dirichlet
边界条件受到高度约束。给定超对称变分 $\delta \lambda= \half
\sum_{IJ}\Gamma^{IJ}F_{IJ}\epsilon$，我们希望选取边界值使得 % 式 (12.15)
\setcounter{equation}{14}
\begin{equation}\label{eq:humbo}
\sum_{I,J\not=
0,1,4}\G^{IJ}F_{IJ}\epsilon=0.
\end{equation}
 如果这对所有 $\epsilon$ 都要成立，我们就得置 $F_{IJ}=0$
（$I,J\not= 0,1,4$）。由于 $\epsilon$ 被约束为
$\Gamma_{2367}\epsilon=T\epsilon=\epsilon$，对 $F$ 的约束就没那么强。
取 $F$ 在由 $x^2,x^3,x^6$ 与 $x^7$ 生成的子空间内自对偶即可。于是
我们必须施加 Hitchin 方程，用这里的记号写出即
% 式 (12.16)
\setcounter{equation}{15}
\begin{equation}\label{eq:ibob}
\begin{aligned}F_{23}-F_{67}& =0\\ F_{27}-F_{36}& = 0 \\
 F_{26}-F_{73}& = 0 . \\\end{aligned}
\end{equation}
  于是规范场的边界值定义了 Higgs 丛模空间 $\MH$ 中的一点 $p$。一旦
施加 \eqref{eq:ibob}，则只要再要求：除非 $I$ 与 $J$ 都属于集合
${\EUR}=\{2,3,6,7\}$，否则 $F_{IJ}|=0$，\eqref{eq:humbo}\ 就成立。
特别地，对 $I\in {\EUR}$、$J\notin {\EUR}$ 有 $F_{IJ}|=0$，这蕴含
$A_J|$（$J=5,8,9$）在（局部坐标为 $x^2,x^3$ 的黎曼面 $C$ 上）协变常数，
并与 Higgs 场交换，而后者的分量为 $\phi_2=A_6 $ 与 $\phi_3=A_7$。
若 $p$ 是 $\MH$ 中对应于不可约 Higgs 丛的光滑点，这些条件蕴含对
$J=5,8,9$ 有 $A_J|=0$。这又进一步蕴含对 $I,J\notin {\EUR}$ 有
$F_{IJ}|=0$。所以对不可约 Higgs 丛，（这一类型的）边界条件被唯一确定。
若 $p$ 对应可约 Higgs 丛，则一般只能得知 $A_5,$ $A_8$ 与 $A_9$ 生成该
Higgs 丛的对称性，且彼此交换。

到目前为止得到的边界条件对应于 $\MH$ 上支撑在点 $p$ 的零膜。这是
$(B,B,B)$-膜的基本类型。正如所料，若 $p$ 对应 $\MH$ 的光滑点，这样的
膜是唯一的。

\par\medskip\emph{零膜的另一种构造}

有点令人意外的是，$\MH$ 上的零膜还有另一种规范场论构造。只要身处
$\MH$ 的光滑点，它在低能理论中似乎给出同一个膜；但对支撑在 $\MH$ 的
奇点处的零膜，它可能给出新的东西。

让我们把 $T$ 换成 $\tilde T = T
\Gamma_{2367}=\Gamma_{023467}$。此时边界条件保持的超对称，其生成元
满足 $\tilde T\epsilon=\epsilon$。于是四维规范场论中的边界条件将保持
一组与先前不同的八个超荷。然而，低能理论中未破缺的超对称是那些满足
组合条件 $T\epsilon=\Gamma_{2367}\epsilon=\epsilon$ 的超对称。这些条件
在把 $T$ 换成 $\tilde T$ 下不变，所以边界条件的改变不会改变低能理论
中未破缺的超对称，而这一构造将给出一个新的 $(B,B,B)$-膜。

我们把费米子上的边界条件取为
$\tilde T\lambda\vert = -\lambda\vert$。要求这一条件在超对称下不变，
我们需要 $\{\Gamma^{IJ}F_{IJ},\tilde T\}\epsilon=0$。让我们再假定
$\Gamma_{2367}\eps=\eps$。于是超对称条件在如下两式成立时满足：% 式 (12.17)
\setcounter{equation}{16}
\begin{equation}\label{eq:kipo}
\sum_{I,J\in
\{0,2,3,4,6,7\}}\Gamma^{IJ}F_{IJ}\vert\epsilon=0
\end{equation}
 以及
% 式 (12.18)
\setcounter{equation}{17}
\begin{equation}\label{eq:kipok}
\sum_{I,J\in
\{1,5,8,9\}}\Gamma^{IJ}F_{IJ}\vert\epsilon=0.
\end{equation}
 若 $A_I$ 满足 Dirichlet 边界条件、$A_4\vert=0$，并且 $A_I$
（$I\in \{2,3,6,7\}$）的边界值被约束为满足 Hitchin 方程
\eqref{eq:ibob}，则 \eqref{eq:kipo}\ 成立。

注意，若 $\eps$ 不满足 $\Gamma_{2367}\eps=\eps$，超对称约束一般不能
满足。这意味着即使在平坦时空中，这个 $(B,B,B)$-膜也只保持四个超荷。

为满足 \eqref{eq:kipok}，我们\emph{不能}要求 $F_{IJ}|=0$
（$I,J\in \{1,5,8,9\}$），因为这些条件是超定的。的确，
$F_{15}|=F_{18}|=0$ 意味着 $A_5$ 与 $A_8$ 满足 Neumann 边界条件、可以
在边界上自由涨落；既然如此，我们就不能再要求 $F_{58}|=0$。然而，
条件 $\tilde T\epsilon=\epsilon$ 连同十维超 Yang-Mills 理论的手征条件
$\hat\Gamma\epsilon=\epsilon$（这里
$\hat\Gamma=\Gamma_0\Gamma_1\cdots \Gamma_9$）蕴含
$\Gamma_{1589}\epsilon=-\epsilon$。在这一对 $\epsilon$ 的约束下，
\eqref{eq:kipok}\ 等价于 $1589$ 平面内边界值上的一个反自对偶条件：
% 式 (12.19)
\setcounter{equation}{18}
\begin{equation}\label{eq:ubb}
\begin{aligned}F_{15}\vert+F_{89}\vert & = 0 \\
                  F_{18}\vert+F_{95}\vert & = 0 \\
                  F_{19}\vert+F_{58}\vert & = 0. \\\end{aligned}
\end{equation}

这是 $A_5,$ $A_8$ 与 $A_9$ 上 Neumann 边界条件的一个修改版本，并且
\emph{不}是超定的。

由于 $A_2,A_3, A_6$ 与 $A_7$ 满足由 Hitchin 方程的某个特定解所确定的
Dirichlet 边界条件，我们这样得到的是支撑在 $\MH$ 中点 $p$ 的零膜。它
与先前零膜构造的区别在于 $A_4,A_5,A_8$ 与 $A_9$ 上的边界条件，以及
在平坦时空中前者保持的超对称只有后者的一半这一点。若 $p$ 是 $\MH$ 中
对应于不可约 Higgs 丛的光滑点，则无论微观边界条件如何，场
$A_4,A_5,A_8$ 与 $A_9$ 都没有零模，并在低能理论中消失。因此，只要 $p$
是 $\MH$ 中的光滑点，零膜的新构造就应与旧构造等价。若 $p$ 对应可约
Higgs 丛，且其稳定化子是 $G$ 的正秩子群，则 $A_4$, $A_5$, $A_8$ 与
$A_9$ 确实有零模——它们生成该 Higgs 丛的对称性——此时没有理由指望
两个零膜等价。

把费米子的边界条件取为 $\tilde T\lambda\vert =
\lambda\vert$，还能得到又一个 $(B,B,B)$-膜。结果表明它对应一个带平坦
规范场的充满空间的膜。由于这个膜在本文中不起主要作用，我们略去细节。

\begin{yizhu}[易错点]
按：$1/4$ BPS 与 $1/2$ BPS 的称呼容易混淆，区别只在比较的基准：四维
${\EUN}=4$ 超 Yang-Mills 有 16 个整体超荷，边界条件保持 4 个，故在
规范场论语言中称 $1/4$ BPS 膜；而紧化后的 sigma 模型只有 8 个超荷，
保持 4 个即占一半，故在同一篇文章的二维语言中称 $1/2$ BPS 膜。读到
后文时务必注意所指的是哪一套记账方式。
\end{yizhu}

\begin{yizhu}[物理图像]
按：本节的物理判据只有一条：边界条件保超对称，当且仅当超流的法向分量
在边界上为零，即 $\bar\epsilon J_1=0$，从而超流通量不会「漏」到边界外。
实际操作顺序是：先对费米子提出简单边条件（保证
${\bar\lambda}\Gamma_1\delta\lambda=0$），由超对称变分推出玻色子边
条件，再回头验证 $J_1=0$。这也解释了为什么边界数据最终落在 Hitchin
方程上：只有满足 Hitchin 方程的边界值，才能让 $\delta\lambda$ 在未破缺
的 $\epsilon$ 上为零。
\end{yizhu}

\begin{yizhu}[易错点]
按：Neumann 与 Dirichlet 在扭化理论中的对应关系与通常场论相反地绕：
$\delta A_I=0$ 是 Dirichlet（固定 $A_I$ 的边界值），而 $F_{I1}=0$ 才是
Neumann 的规范不变版本（$A_1=0$ 规范下即 $A_I$ 法向导数为零）。$A_0+A_4$
是「零」方向：由于号差中含负号（\eqref{eq:ofof} 中
$g^{IJ}={\rm diag}(-1,1,1,\dots,1)$），双曲边界条件可以同时把
$A_0+A_4$ 的值与法向导数置零，这不违反双曲性。
\end{yizhu}

\begin{yizhu}[含糊点拨]
按：本片有两处引用了第 5 节的式子：式 (12.1) 那句「正如式
\eqref{eq:ytytag} 那样」中的 \eqref{eq:ytytag} 指向
skel/11-s05.tex 中带标签 eq:ytytag 的式子，即
$\Gamma_{2367}\epsilon=\epsilon$（按：此处应读作式 (5.3)）；该标签所在
公式在骨架中编号未落定（生成时留有 WARN 注释），因此在全书标签统一
修正前可能显示为 ??，读者按第~\ref{sec:reviewsym}~节的相关讨论理解即可。
\end{yizhu}

\begin{yizhu}[物理图像]
按：$(B,B,B)$ 型膜的规范场论实现可以一句话概括：边界值是 Hitchin 方程
的一个解，因而定义了 $\MH$ 中的一点 $p$，膜就支撑在 $p$ 上。对不可约
Higgs 丛边界条件唯一确定；对可约 Higgs 丛，$A_5,A_8,A_9$ 只被确定为
生成 Higgs 丛的对称性（零模），这正是不唯一性的来源，也解释了为什么
在奇异点处两种构造可能给出不同的膜。
\end{yizhu}

\begin{yizhu}[原文疑误]
按：本片有两处原文问题，均照录并说明。(1) 第 49 行 "unfelicitous" 应为
"infelicitous"（不妥当的）。(2) 倒数第二段「在平坦时空中前者保持的
超对称只有后者的一半」一句，若按英文 former/latter 的字面指代
（former = 先前构造，latter = 新构造），则与前文「新构造即使在平坦
时空中也只保持四个超荷」的论述相反：先前的 $T=\G_{04}$ 构造在平坦时空
中保八个超荷，新构造只保四个。故此处 former/latter 疑为颠倒，或作者
笔下的指代与本段开头「it 与 the previous construction」的顺序不一致；
译文保留原文语序。
\end{yizhu}
